% GENERATED FILE. Edit canonical lessons, then run npm run generate:books.
% canonical-source-hash: fnv1a64:82181cd5b34f393b
% canonical-lessons: TE-C24-ratri, TE-S127-letter-ja

\chapter{Night}
\label{ch:te-night}
\hlchaptermodality{24}

\begin{chapteropening}
\textbf{By the end of this chapter:} \emph{I can say \textquotedblleft{}night\textquotedblright{} in Telugu and recognise it inside the midnight word I already knew.}

Say \te{రాత్రి}, find it again inside \te{అర్ధరాత్రి}, and note the three native words — rēyi, māpu, irulu — it pushed aside.
\end{chapteropening}

\section[rātri]{\te{రాత్రి} (rātri) — \textquotedblleft{}night,\textquotedblright{} this time matching the expected pattern}

\label{lesson:TE-C24-ratri}



\begin{quote}
Last lesson's \emph{rōju} broke every expectation. This one is more straightforward — but still worth checking honestly, not assuming.
\end{quote}

\begin{cousinweb}
\textbf{\te{రాత్రి}} (\textbf{rātri}) — \textquotedblleft{}\textbf{night}\textquotedblright{} — is a Sanskrit \textbf{tatsama} word, already hiding inside \textbf{\te{అర్ధరాత్రి}} (\emph{ardharātri}, \textquotedblleft{}midnight,\textquotedblright{} Chapter 17). Unlike \emph{rōju}'s Persian surprise, \emph{rātri} really \textbf{is} Telugu's ordinary, everyday word for \textquotedblleft{}night\textquotedblright{} — matching Kannada's own pattern this time, not breaking it.
\end{cousinweb}

\begin{cousinweb}
Telugu \textbf{does} have native words here — but be honest: it's a \textbf{messier} picture than Kannada's clean two-word story, not a tidy mirror of it.

\begin{itemize}
\raggedright
  \item \textbf{\te{రేయి}} (\textbf{rēyi}) — attested in older Telugu literature; most sources treat it as inherited Dravidian, though some grammars instead link it to \emph{rātri} itself, so its native status is less than fully settled.
  \item \textbf{\te{మాపు}} (\textbf{māpu}) — its \textbf{primary} sense is actually \textquotedblleft{}\textbf{evening, dusk}\textquotedblright{} (as in \textbf{\te{రేపుమాపు}}, \textquotedblleft{}morning and evening\textquotedblright{}), with \textquotedblleft{}night\textquotedblright{} only a secondary, extended sense — not a clean synonym for \emph{rātri}.
  \item \textbf{\te{ఇరులు}} (\textbf{irulu}) — a genuine cognate of Kannada's \textbf{\ka{ಇರುಳು}} (\emph{iruḷu}), but its primary modern sense is \textquotedblleft{}\textbf{darkness},\textquotedblright{} archaic even in that use.
\end{itemize}

All three have been effectively pushed aside by the Sanskrit loan — but unlike Kannada's single native competitor, Telugu's \textquotedblleft{}native\textquotedblright{} side is split across three words in three different, overlapping senses.
\end{cousinweb}

\subsection*{Your turn}
Say these aloud:

\begin{itemize}
\raggedright
  \item \textquotedblleft{}rātri\textquotedblright{} — \textquotedblleft{}night,\textquotedblright{} already met inside ardharātri
  \item rēyi, māpu, irulu — three native words, three overlapping senses, all pushed aside
  \item the contrast with last lesson — day broke the pattern, night matches it
\end{itemize}

\begin{tcolorbox}[breakable,colback=teal!4,colframe=teal!35!black,title={Before you move on}]
Where have you already met \textbf{\te{రాత్రి}} before this lesson? (\textbf{\te{అర్ధరాత్రి}}, \emph{ardharātri}, \textquotedblleft{}midnight.\textquotedblright{}) Does \emph{rātri} follow the same \textquotedblleft{}Sanskrit-as-everyday\textquotedblright{} pattern as Kannada's \emph{rātri}, or does it break the pattern the way \emph{rōju} did for \textquotedblleft{}day\textquotedblright{}? (\textbf{It matches} — \emph{rātri} really is Telugu's everyday word.) Does Telugu's native-word story cleanly mirror Kannada's single-word \emph{iruḷu}? (\textbf{No} — it's messier: three words, \textbf{\te{రేయి}}/\textbf{\te{మాపు}}/\textbf{\te{ఇరులు}}, in three overlapping senses, not one clean replacement.)
\end{tcolorbox}
\section[ja]{\te{జ} — one character, met inside words you already say}

\label{lesson:TE-S127-letter-ja}



\begin{quote}
Before the new one: \te{ణ} — what does it do? And one from further back: \te{ళ}?

One character this time — and you have been saying it for pages without knowing which mark on the page it was.
\end{quote}

\subsection*{Script you'll notice: \te{జ}}
\textbf{\te{జ}} — \emph{ja}.

It is a \textbf{consonant}, and in this script a consonant is never bare: it comes with an \emph{a} already in it. So it is not \emph{j}, it is \textbf{ja}.

You already say these, and every one of them has \te{జ} somewhere inside it:

\begin{itemize}
\raggedright
  \item \textbf{\te{రోజు}} \emph{rōju} — day, the word that came in from Persian
  \item \textbf{\te{జ్యేష్ఠం}} \emph{Jyēṣṭhaṁ} — the third of the twelve months
  \item \textbf{\te{ఆశ్వయుజం}} \emph{Āśvayujaṁ} — the seventh
\end{itemize}

\emph{Rōju} is worth a second look. Its first letter is the \textbf{\te{ర}} you traced inside \emph{namaskāram}. Its last is \textbf{\te{జు}} — this new letter carrying the \emph{u} sign you met early on. The vowel sign in the middle is one this book has not stopped on yet, so the word is not wholly yours on the page. Two of its three pieces are, and that is new today.

\subsection*{Writing: \te{జ} — copy what you see}
Put your pen on \te{జ} and follow its line. Copy the shape you can see — slowly, and larger than it is printed.

\begin{quote}
This book does not yet tell you \textbf{where to start the character or which way to travel}. That is a real thing, taught with real variation from school to school, and it is not written down here until it can be written down with a source. Copying what is in front of you needs no such source, and it is how the shape gets into your hand in the meantime.
\end{quote}

\subsection*{Your turn}
\begin{itemize}
\raggedright
  \item \emph{Look:} at these words, and find \te{జ} in the ones that have it
\end{itemize}

\begin{quote}
\te{రోజు}  ·  \te{నమస్కారం}
\end{quote}

\begin{itemize}
\raggedright
  \item \emph{Trace it:} \te{జ} three times, saying \emph{ja} as you finish each one
  \item \emph{Look:} back at any page of this chapter and find \te{జ} once more
\end{itemize}

\begin{tcolorbox}[breakable,colback=teal!4,colframe=teal!35!black,title={Before you move on}]
Which character is this — \te{జ}? What sound does it carry? (\emph{\textbf{ja}}.) Read \te{రోజు} aloud, letter by letter, and then as one word.

One more, from much earlier: \te{◌ె} — what does it do?
\end{tcolorbox}
